% paper/sections/08e_ppe_bc.tex
%
% §8.5  境界条件と数値安定性
%
\subsection{境界条件と数値安定性}
\label{sec:ppe_bc_stability}
\begin{itemize}
    \item \textbf{固体壁・流入出：} $\partial p/\partial n = 0$（ノイマン条件）．CCD ブロック方程式への組み込み手順（境界行を $p'_0=0$ で置き換え，変密度項の消滅）は §~\ref{sec:ccd_neumann_bc} を参照．
    \item \textbf{自由表面：} $p = p_{\mathrm{atm}}$（ディリクレ条件）
    \item \textbf{全周ノイマンの場合：} 解の一意性を保つため，1点の圧力を固定するか平均値を0にする拘束を加える（§~\ref{sec:pinpoint}参照）．
    \item \textbf{周期境界：} 全辺を周期とする場合は FVM 行列に折り返しフェイスおよびゴーストノード拘束を組み込む（§~\ref{sec:fvm_periodic} 参照）．零空間対策は壁面 BC と同一（1点ピン条件）．
\end{itemize}

%% Neumann BC 実装検証セット（無負荷試験・製造解収束テスト・零モード確認）は
%% 付録~\ref{sec:ccd_neumann_unit_test} に移動（CHK-086）．
%% FVM PPE 周期境界条件対応（折り返しフェイス・ゴーストノード拘束・ILU 充填）は
%% 付録~\ref{sec:fvm_periodic} に移動（CHK-086）．

\subsubsection{高密度比における条件数と数値安定性}
\label{sec:ppe_condition_number}

気液界面を挟む密度比 $\rho_l/\rho_g \approx 10^3$（水/空気相当）は
PPE 演算子の条件数に直接影響する．
可変係数演算子 $\mathcal{L}_h$ の係数 $a_{i+1/2} = 2/(\rho_i+\rho_{i+1})$ は，
液相（$\rho \approx \rho_l$）と気相（$\rho \approx \rho_g$）の格子フェイスで
$O(\rho_l/\rho_g)$ の比率を持つため，行列の条件数は
\begin{equation}
  \kappa(\mathcal{L}_h) = O\!\left(\frac{\rho_l}{\rho_g}\cdot\frac{1}{h^2}\right)
  \approx O\!\left(\frac{10^3}{h^2}\right)
  \label{eq:ppe_cond_number}
\end{equation}
と見積もられる．

\noindent\textbf{安定化戦略：}
\begin{itemize}
  \item \textbf{調和平均フェイス係数（式~\eqref{eq:af_exact}）：}
        単純算術平均（$(\rho_L+\rho_R)/2$ の逆数）と比較して，
        調和平均 $2/(\rho_L+\rho_R)$ は界面フラックスに物理的根拠を持ち，
        条件数の過大評価を防ぐ（付録~\ref{app:fvm_face_coeff}）．
  \item \textbf{Extension PDE による界面越し圧力延長（§~\ref{sec:extension_pde}）：}
        PPE 求解前に圧力場を Extension PDE で界面越しに滑らかに延長することで，
        CCD ステンシル $(i-2,\ldots,i+2)$ が界面直近まで $C^\infty$ に近い場を参照でき，
        Gibbs 振動が除去されて反復ソルバーの収束が安定化される．
  \item \textbf{ゲージ固定点の対称性（KL-11；§~\ref{sec:pinpoint}）：}
        ピン固定点を中心ノード $(N/2, N/2)$ に設定することで（コーナーピン禁止），
        高密度比条件下でもグローバル非対称性 $O(\|\mathrm{rhs}\|)$ を回避する
        （ASM-008 対称性修正，2026-03-22 適用済み）．
  \item \textbf{LGMRES の頑健性：}
        CCD 境界スキームによる非対称性の強い行列でも LGMRES（拡張 Krylov 部分空間法）は
        標準 GMRES より頑健に収束する（付録~\ref{app:ccd_lu_direct}）．
        収束しない場合は最良反復解を返す．
\end{itemize}

\subsubsection{全周 Neumann 境界における零モード対策}
\label{sec:ppe_zero_mode}

全境界が $\partial p/\partial n = 0$（ノイマン条件）の場合，
PPE の係数行列は特異（零固有値あり）であり，解 $p$ は定数の自由度を持つ．
この自由度を取り除かないと反復ソルバーが発散する．

\begin{enumerate}
  \item \textbf{1点固定（ピン条件，推奨）：}
        セル $(i_0, j_0)$ の圧力を $p_{i_0,j_0} = 0$ に固定し，その行を単位行ベクトルに置き換える．
        Matrix-free 反復では各ステップ後に強制上書きする：
        \begin{equation}
          (\delta p)^m_{i_0, j_0} \leftarrow 0
          \qquad \text{（全反復ステップで更新直後に適用）}
          \label{eq:pin_overwrite}
        \end{equation}
  \item \textbf{平均値拘束：} $\sum_{i,j} p_{i,j} = 0$ の線形拘束を追加する（精度的に優位だが実装が複雑）．
\end{enumerate}

\subsubsection{圧力固定点（ピン点）の選択原則}
\label{sec:pinpoint}

\noindent\textbf{原則1：固定点を途中で変更しない．}
ピン点を変更すると，圧力の基準値が不連続に変化し圧力勾配に突発的誤差が生じる．

\noindent\textbf{原則2：固定点は純粋気相または純粋液相の領域に置く．}
界面が固定点を通過すると，ピン条件により正しい PPE 方程式が反映されず収束が崩れる．
4倍対称な設定では中心セル $(N/2,N/2)$ が数値対称性を保持する
（KL-11；§~\ref{sec:ppe_gauge_neumann} 参照）．

\noindent\textbf{原則3：圧力のゼロ点は相対量である．}
$p_{i_0,j_0} = 0$ は数学的拘束であり，速度補正に必要な勾配 $\bnabla p$ はスカラーシフトに依存しない．
